\documentclass[11pt,a4paper]{article}
\usepackage[margin=20mm,top=27mm,bottom=20mm,headheight=25pt,footskip=8mm]{geometry}
\usepackage[utf8]{inputenc}
\usepackage[T1]{fontenc}
\usepackage[brazil]{babel}
\usepackage{helvet,amsmath,amssymb,booktabs,array,tabularx,enumitem,graphicx,xcolor,tcolorbox,fancyhdr,hyperref,tikz}
\usetikzlibrary{arrows.meta,positioning}
\renewcommand{\familydefault}{\sfdefault}
\definecolor{azul}{RGB}{0,83,134}
\definecolor{ouro}{RGB}{206,177,68}
\hypersetup{pdflang=pt-BR,colorlinks=true,urlcolor=azul,linkcolor=azul,pdfauthor={Luan F. Nogueira e Luiz F. de Melo Soares}}
\setlength{\parindent}{0pt}
\setlength{\parskip}{5pt}
\setlist{nosep,leftmargin=*}
\pagestyle{fancy}\fancyhf{}
\fancyhead[L]{\small\color{azul}DO TRANSISTOR AO CHIP}
\fancyhead[R]{\includegraphics[height=7mm]{../slides/libs/emblemufc.pdf}}
\fancyfoot[L]{\scriptsize Luan F. Nogueira • Luiz F. de Melo Soares}
\fancyfoot[R]{\scriptsize\thepage}
\renewcommand{\headrulewidth}{.4pt}
\setlength{\headheight}{25pt}
\newcommand{\titulo}[2]{{\LARGE\color{azul}\bfseries #1}\par{\large #2}\par\vspace{2mm}}
\newcommand{\caixa}[2]{\begin{tcolorbox}[colback=azul!4,colframe=azul,title={#1},fonttitle=\bfseries,boxrule=.5pt,arc=1mm]#2\end{tcolorbox}}
\newcommand{\VDD}{V_{\mathrm{DD}}}

\hypersetup{pdftitle={Gabarito comentado do minicurso Do transistor ao chip}}
\begin{document}
\titulo{Gabarito comentado}{Atividades das aulas 1 e 2}
\textbf{1 · Materiais e portadores.} a) Tipo P: lacunas; tipo N: elétrons. b) Os íons dos dopantes e os portadores móveis se compensam aproximadamente no volume neutro. c) Íons fixos dos dopantes: aceitadores negativos no lado P e doadores positivos no lado N; há poucos portadores móveis na depleção. d) Polarização direta: P em potencial maior que N.

\textbf{2 · Diferenças de tensão.}
\begin{center}\begin{tabular}{llll}\toprule Caso & Cálculo & Comparação & Estado ideal\\\midrule
NMOS, G=0; S=0 & $V_{GS}=0-0=0$ V & $0<0{,}4$ V & Cortado\\
PMOS, G=0,2; S=1,2 & $V_{SG}=1{,}2-0{,}2=1{,}0$ V & $1{,}0>0{,}4$ V & Ligado\\
NMOS, G=1,2; S=1,0 & $V_{GS}=1{,}2-1{,}0=0{,}2$ V & $0{,}2<0{,}4$ V & Cortado\\
PMOS, G=1,2; S=1,2 & $V_{SG}=1{,}2-1{,}2=0$ V & $0<0{,}4$ V & Cortado\\\bottomrule\end{tabular}\end{center}
A terceira linha mostra por que olhar apenas a tensão da porta pode levar ao erro. Aqui se idealiza um limiar constante; o dispositivo real também depende de corpo e condições de operação.

\textbf{3 · Inversor e nó flutuante.} a) O PMOS conduz e carrega Y até próximo de $\VDD$: saída lógica 1. O NMOS está cortado. b) Sem PMOS e sem outro elemento de carga, não há caminho garantido para recarregar a saída: ela fica flutuante quando o NMOS corta. c) Sem caminho ativo, a tensão pode depender da carga anterior, das capacitâncias e das fugas; circuito aberto não equivale a terra.

\textbf{4 · Portas de duas entradas.} NAND: PMOS em paralelo entre alimentação e saída; NMOS em série entre saída e GND. Para A=1 e B=0, o PMOS B permite a carga, e o NMOS B interrompe a descida. NOR: PMOS em série e NMOS em paralelo.
\begin{center}\begin{tabular}{cccccc}\toprule A & B & NAND: subida & NAND: Y & NOR: descida & NOR: Y\\\midrule
0&0&Sim&1&Não&1\\0&1&Sim&1&Sim&0\\1&0&Sim&1&Sim&0\\1&1&Não&0&Sim&0\\\bottomrule\end{tabular}\end{center}
Em todos os desenhos, as portas identificadas por A recebem a mesma entrada A; o mesmo vale para B. Os corpos NMOS ficam em GND e os corpos PMOS em $\VDD$ no modelo introdutório.

\textbf{5 · NAND3.} Três NMOS em série e três PMOS em paralelo: seis transistores. A saída é zero somente em 111. Em qualquer outra combinação, a cadeia NMOS é interrompida e pelo menos um PMOS fornece caminho até a alimentação.

\textbf{6 · Justificativas.} a) A série define a condição AB para descarregar a saída; portanto $Y=\overline{AB}$. b) Com 01, a série NMOS tem A aberto e a série PMOS tem B aberto: nenhuma rede fixa a saída. c) Unir as duas entradas: $N(A,A)=\overline{AA}=\overline A$. d) Carregar e descarregar capacitâncias envolve energia. Há ainda fuga e possível corrente direta transitória. Baixo consumo estático não significa energia de comutação nula.
\newpage
\titulo{Extensão opcional · AOI21}{Combinar as regras de série e paralelo}
\[Y=\overline{AB+C}\]
\textbf{Rede NMOS:} A em série com B; esse ramo em paralelo com C.\\
\textbf{Rede PMOS:} A em paralelo com B; esse conjunto em série com C.\\
\textbf{Contagem:} três NMOS e três PMOS, seis transistores.

\begin{center}\begin{tabular}{cccccc}\toprule A & B & C & Descida ativa? & Subida ativa? & Y\\\midrule
0&0&0&Não&Sim&1\\0&0&1&Sim&Não&0\\0&1&0&Não&Sim&1\\0&1&1&Sim&Não&0\\
1&0&0&Não&Sim&1\\1&0&1&Sim&Não&0\\1&1&0&Sim&Não&0\\1&1&1&Sim&Não&0\\\bottomrule\end{tabular}\end{center}
\caixa{Como avaliar o raciocínio}{A resposta deve relacionar a expressão, a rede e cada caminho de condução. Verifique se o participante distingue a condição de condução da rede NMOS do valor lógico na saída.}
\textbf{Universalidade · consulta adicional.} Definindo $N(A,B)=\overline{AB}$ e $R(A,B)=\overline{A+B}$:
\begin{align*}
\overline A&=N(A,A)=R(A,A)\\
AB&=N\bigl(N(A,B),N(A,B)\bigr)=R\bigl(R(A,A),R(B,B)\bigr)\\
A+B&=N\bigl(N(A,A),N(B,B)\bigr)=R\bigl(R(A,B),R(A,B)\bigr).
\end{align*}
\textbf{XOR com quatro NANDs.} Com $T=N(A,B)$, $U=N(A,T)$, $V=N(B,T)$ e $Y=N(U,V)$, os valores de Y para 00, 01, 10, 11 são 0, 1, 1, 0. A implementação usa 16 transistores em quatro NAND2 de quatro transistores; não é apresentada como contagem mínima.

\small\textbf{Referências:} Fathi Salem, \textit{ECE 410: VLSI Design}, e J.-M. Dutertre, \textit{Digital Design (Introduction to)}, PDFs fornecidos. Síntese didática e gabarito pelos facilitadores. Todas as conclusões de caminhos consideram entradas binárias estáveis e o modelo ideal de CMOS estático complementar.
\end{document}
